% GENERATED FILE. Edit canonical lessons, then run npm run generate:books.
% canonical-source-hash: fnv1a64:941ec49e5037581a
% canonical-lessons: RU-C56-mountain, RU-C56-forest, RU-C56-river, RU-C56-lake, RU-C56-sea

\chapter{Mountain, Forest, River, Lake, Sea}
\label{ch:ru-mountain-forest-river-lake-sea}
\hlchaptermodality{56}

\begin{chapteropening}
\textbf{By the end of this chapter:} \emph{I can name a mountain, a forest, a river, a lake and the sea.}

Complete the last lesson of chapter 56: I can name a mountain, a forest, a river, a lake and the sea.
\end{chapteropening}

\section[\texorpdfstring{\ru{гора} (gorá)}{гора (gorá)}]{\ru{гора} (gorá) — a mountain}

\label{lesson:RU-C56-mountain}



\begin{quote}
Before the new one: say the Russian for grass, then the Russian for a village.
\end{quote}

\subsection*{You'll want to know: \ru{гора}}
\textbf{\ru{гора}} (\emph{gorá}) — \textquotedblleft{}a mountain\textquotedblright{}.

\begin{cousinweb}
\textbf{\ru{гора}} $\leftarrow$ Proto-Slavic \emph{*gora} — inherited, related to Sanskrit \emph{girí-}, \textquotedblleft{}mountain\textquotedblright{}, Avestan \emph{gairi-} and Lithuanian \emph{girià}, \textquotedblleft{}forest\textquotedblright{} (Vasmer). Bulgarian \emph{\ru{гора}} still means \textquotedblleft{}forest\textquotedblright{}.
\end{cousinweb}

\begin{grammarlens}[title={What you've built}]
One more word for this chapter.
\end{grammarlens}

\subsection*{Your turn}
\begin{itemize}
\raggedright
  \item \emph{Say it:} \emph{gorá}
  \item \emph{Say it:} \emph{gorá}, once more
  \item \emph{Say it:} say \emph{gorá}
  \item \emph{Recall:} say the Russian for grass, then the Russian for a village, then say \emph{gorá} again
\end{itemize}

\begin{tcolorbox}[breakable,colback=teal!4,colframe=teal!35!black,title={Before you move on}]
What does \ru{гора} mean? (\textquotedblleft{}A mountain\textquotedblright{}.) Say it once more.
\end{tcolorbox}
\section[\texorpdfstring{\ru{лес} (les)}{лес (les)}]{\ru{лес} (les) — a forest}

\label{lesson:RU-C56-forest}



\begin{quote}
Before the new one: say the Russian for a village, then the Russian for a mountain.
\end{quote}

\subsection*{You'll want to know: \ru{лес}}
\textbf{\ru{лес}} (\emph{les}) — \textquotedblleft{}a forest\textquotedblright{}.

\begin{cousinweb}
\textbf{\ru{лес}} $\leftarrow$ Proto-Slavic \emph{*lěsŭ} — inherited. Its deeper origin is not certain: Vasmer relates it to Old English \emph{læs}, \textquotedblleft{}pasture\textquotedblright{}, and English dialect \emph{leasow}, and finds a link with Lithuanian \emph{laĩškas}, \textquotedblleft{}leaf\textquotedblright{}, doubtful, while Trubachev, in a note, thinks it the most likely.
\end{cousinweb}

\begin{grammarlens}[title={What you've built}]
One more word for this chapter.
\end{grammarlens}

\subsection*{Your turn}
\begin{itemize}
\raggedright
  \item \emph{Say it:} \emph{les}
  \item \emph{Say it:} \emph{les}, once more
  \item \emph{Say it:} say \emph{les}
  \item \emph{Recall:} say the Russian for a village, then the Russian for a mountain, then say \emph{les} again
\end{itemize}

\begin{tcolorbox}[breakable,colback=teal!4,colframe=teal!35!black,title={Before you move on}]
What does \ru{лес} mean? (\textquotedblleft{}A forest\textquotedblright{}.) Say it once more.
\end{tcolorbox}
\section[\texorpdfstring{\ru{река} (reká)}{река (reká)}]{\ru{река} (reká) — a river}

\label{lesson:RU-C56-river}



\begin{quote}
Before the new one: say the Russian for a mountain, then the Russian for a forest.
\end{quote}

\subsection*{You'll want to know: \ru{река}}
\textbf{\ru{река}} (\emph{reká}) — \textquotedblleft{}a river\textquotedblright{}.

\begin{cousinweb}
\textbf{\ru{река}} $\leftarrow$ Proto-Slavic \emph{*rěka} — inherited, from a root \textquotedblleft{}to flow\textquotedblright{}. Vasmer lists Sanskrit \emph{rīyate}, \textquotedblleft{}begins to flow\textquotedblright{}, Latin \emph{rīvus}, \textquotedblleft{}a stream\textquotedblright{}, and Gaulish \emph{Rēnos}, the Rhine.
\end{cousinweb}

\begin{grammarlens}[title={What you've built}]
One more word for this chapter.
\end{grammarlens}

\subsection*{Your turn}
\begin{itemize}
\raggedright
  \item \emph{Say it:} \emph{reká}
  \item \emph{Say it:} \emph{reká}, once more
  \item \emph{Say it:} say \emph{reká}
  \item \emph{Recall:} say the Russian for a mountain, then the Russian for a forest, then say \emph{reká} again
\end{itemize}

\begin{tcolorbox}[breakable,colback=teal!4,colframe=teal!35!black,title={Before you move on}]
What does \ru{река} mean? (\textquotedblleft{}A river\textquotedblright{}.) Say it once more.
\end{tcolorbox}
\section[\texorpdfstring{\ru{озеро} (ózero)}{озеро (ózero)}]{\ru{озеро} (ózero) — a lake}

\label{lesson:RU-C56-lake}



\begin{quote}
Before the new one: say the Russian for a forest, then the Russian for a river.
\end{quote}

\subsection*{You'll want to know: \ru{озеро}}
\textbf{\ru{озеро}} (\emph{ózero}) — \textquotedblleft{}a lake\textquotedblright{}.

\begin{cousinweb}
\textbf{\ru{озеро}} $\leftarrow$ Proto-Slavic \emph{*ezero} — inherited, a cousin of Lithuanian \emph{ẽžeras} and Latvian \emph{ezers}, \textquotedblleft{}lake\textquotedblright{} (Vasmer). East Slavic turned the first \emph{je-} into \emph{\ru{о}-}: compare Old Church Slavonic \emph{jezero}.
\end{cousinweb}

\begin{grammarlens}[title={What you've built}]
One more word for this chapter.
\end{grammarlens}

\subsection*{Your turn}
\begin{itemize}
\raggedright
  \item \emph{Say it:} \emph{ózero}
  \item \emph{Say it:} \emph{ózero}, once more
  \item \emph{Say it:} say \emph{ózero}
  \item \emph{Recall:} say the Russian for a forest, then the Russian for a river, then say \emph{ózero} again
\end{itemize}

\begin{tcolorbox}[breakable,colback=teal!4,colframe=teal!35!black,title={Before you move on}]
What does \ru{озеро} mean? (\textquotedblleft{}A lake\textquotedblright{}.) Say it once more.
\end{tcolorbox}
\section[\texorpdfstring{\ru{море} (móre)}{море (móre)}]{\ru{море} (móre) — the sea}

\label{lesson:RU-C56-sea}



\begin{quote}
Before the new one: say the Russian for a river, then the Russian for a lake.
\end{quote}

\subsection*{You'll want to know: \ru{море}}
\textbf{\ru{море}} (\emph{móre}) — \textquotedblleft{}the sea\textquotedblright{}.

\begin{cousinweb}
\textbf{\ru{море}} $\leftarrow$ Proto-Slavic \emph{*morje} — inherited. Vasmer lists Lithuanian \emph{mãrios}, \textquotedblleft{}lagoon\textquotedblright{}, Gothic \emph{marei}, Old Irish \emph{muir} and Latin \emph{mare}, all \textquotedblleft{}sea\textquotedblright{}: English \emph{mere}, \textquotedblleft{}a lake\textquotedblright{}, is a cousin. He gives \textquotedblleft{}marsh\textquotedblright{} as a possible older sense.
\end{cousinweb}

\begin{grammarlens}[title={What you've built}]
One more word for this chapter.
\end{grammarlens}

\subsection*{Your turn}
\begin{itemize}
\raggedright
  \item \emph{Say it:} \emph{móre}
  \item \emph{Say it:} \emph{móre}, once more
  \item \emph{Say it:} say \emph{móre}
  \item \emph{Recall:} say the Russian for a river, then the Russian for a lake, then say \emph{móre} again
\end{itemize}

\begin{tcolorbox}[breakable,colback=teal!4,colframe=teal!35!black,title={Before you move on}]
What does \ru{море} mean? (\textquotedblleft{}The sea\textquotedblright{}.) Say it once more.
\end{tcolorbox}
